\chapter{АДАПТИВНОЕ МНОГОАГЕНТНОЕ УПРАВЛЕНИЕ}\label{sec:ch4}

Во второй главе была построена доменная и математическая модель системы
автономных мобильных роботов с поддержкой периферийных вычислений. В третьей
главе на основе этой модели был введён VCG-подобный аукционный слой, который
распределяет вычислительные задачи между периферийными узлами при
фиксированных операционных режимах этих узлов.

Однако аукционный слой сам по себе решает только задачу текущего раунда. Он
не определяет, какую часть ресурсов узел должен экспонировать для новых задач,
когда следует резервировать ёмкость под будущие критические запросы, как
реагировать на рост очередей, ухудшение связи или временную деградацию узла.
Эти вопросы относятся к последовательному адаптивному управлению.

Цель настоящей главы состоит в построении обучаемого многоагентного контура,
который выбирает операционные режимы периферийных узлов. В предлагаемой
постановке обучаемыми агентами являются не автономные мобильные роботы, а
периферийные вычислительные узлы. Роботы формируют вычислительные задачи и
заявки для аукционного слоя, но не обучают стратегию искажения заявок или
оценок ценности. Такое разделение необходимо для сохранения экономической
интерпретируемости VCG-подобного механизма.

В главе решаются следующие задачи:
\begin{enumerate}
	\item обосновать необходимость обучаемого адаптивного контура поверх
	раундового аукционного распределения;
	\item формализовать управление периферийными узлами как
	децентрализованный частично наблюдаемый марковский процесс принятия
	решений;
	\item определить глобальное состояние, локальные наблюдения, действия и
	функцию вознаграждения;
	\item обосновать применение парадигмы централизованного обучения и
	децентрализованного исполнения;
	\item обосновать выбор алгоритма QMIX для кооперативного управления
	дискретными режимами периферийных узлов;
	\item зафиксировать границы применимости обучаемого контура.
\end{enumerate}

\section{Почему аукционный слой необходимо дополнить адаптивной политикой}

Аукционный механизм, построенный в предыдущей главе, распределяет задачи в
рамках одного временного раунда. При фиксированных режимах периферийных узлов
он выбирает допустимое назначение, максимизирующее общественное
благосостояние текущего шага. Это свойство важно, но оно не решает всю задачу
управления системой.

В AMR-системе текущее решение влияет на последующее состояние среды. Если
узел в текущем раунде экспонирует слишком большую долю ресурсов, он может
получить высокое значение общественного благосостояния сейчас, но в следующих
раундах столкнуться с ростом очереди, нарушением сроков обработки или
невозможностью принять критическую задачу. Если же узел действует слишком
консервативно, система может отклонять слишком много задач и недоиспользовать
доступные ресурсы.

Кроме того, состояние среды частично наблюдаемо. Отдельный периферийный узел
не знает точных будущих поступлений задач, полной картины загрузки всех
соседних узлов и будущего качества беспроводной связи. Он наблюдает только
локальную загрузку, очередь, оценки задержек до ближайших роботов, результаты
предыдущих назначений и ограниченный набор агрегированных признаков.

Поэтому управление режимами периферийных узлов следует рассматривать как
последовательную многоагентную задачу. Её естественный формализм ---
децентрализованный частично наблюдаемый марковский процесс принятия решений
(Decentralized Partially Observable Markov Decision Process, Dec-POMDP)
\cite{Oliehoek2016DecPOMDP,Bernstein2002Complexity}. Точное решение таких
задач в общем случае вычислительно трудно, поэтому на практике применяются
приближённые методы многоагентного обучения с подкреплением
\cite{Sutton2018RL,Hady2025MARLRAOSurvey,Ismail2025DRLMECSurvey}.

\section{Агенты обучаемой модели}

В предлагаемой постановке каждому периферийному вычислительному узлу
$n_j \in \mathcal{N}$ соответствует обучаемый агент $a_j$. Множество агентов
обозначается как
\begin{equation}
	\mathcal{A}^{agent} = \{a_1,a_2,\ldots,a_k\}.
	\label{eq:ch4_agents}
\end{equation}

Такой выбор агентов обусловлен архитектурой метода AMPLE-AMR. Роботы
генерируют задачи и передают их в механизм распределения. Периферийные узлы,
в свою очередь, управляют тем, какую часть ресурсов открыть для новых задач,
какие классы задач обслуживать приоритетно и когда переходить в режим защиты
от перегрузки. Следовательно, именно периферийные узлы являются естественными
субъектами адаптивного управления.

Данный выбор принципиально отличается от постановки, в которой обучаются
стратегии заявок роботов. Если обучаемая политика изменяет ставки или оценки
ценности задач, то VCG-подобный аукционный слой перестаёт быть механизмом
правдивого раскрытия предпочтений и превращается в повторяющуюся игру
стратегического поведения. В настоящей работе эта ситуация исключается:
роботы не обучают bidding policy, а периферийные узлы обучают только политику
выбора операционных режимов.

\section{Формализация в виде Dec-POMDP}

Задача адаптивного управления периферийными узлами формализуется в виде децентрализованного частично наблюдаемого марковского процесса принятия решений (Dec-POMDP)~\cite{Oliehoek2016DecPOMDP,Amato2024CoopOntro}:
\begin{equation}
	\mathcal{M} =
	\langle
	\mathcal{A}^{agent},
	\mathcal{S},
	\{\mathcal{O}_j\}_{j=1}^{k},
	\{\mathcal{K}_j\}_{j=1}^{k},
	P,
	R,
	\Omega,
	\gamma
	\rangle,
	\label{eq:ch4_decpomdp}
\end{equation}
где:
\begin{itemize}
	\item $\mathcal{A}^{agent}$ --- множество обучаемых агентов, соответствующих
	периферийным узлам;
	\item $\mathcal{S}$ --- множество глобальных состояний системы;
	\item $\mathcal{O}_j$ --- пространство локальных наблюдений агента $a_j$;
	\item $\mathcal{K}_j$ --- множество операционных режимов узла $n_j$;
	\item $P(s_{t+1}\mid s_t,\kappa_t)$ --- вероятностная функция переходов (описывает, как состояние системы меняется в результате действий всех агентов);
	\item $R(s_t,\kappa_t,s_{t+1})$ --- командная функция вознаграждения;
	\item $\Omega(o^1_t,\ldots,o^k_t\mid s_t)$ --- функция наблюдений (описывает, какое локальное наблюдение получит агент после перехода в новое состояние);
	\item $\gamma \in [0,1)$ --- коэффициент дисконтирования (параметр, определяющий вес будущих вознаграждений).
\end{itemize}

Совместное действие агентов в момент времени $t$ задаётся вектором
\begin{equation}
	\kappa_t =\bigl(\kappa_1(t),\ldots,\kappa_k(t)\bigr),
	\label{eq:ch4_global_action}
\end{equation}
где $\kappa_j(t)\in\mathcal{K}_j$ --- операционный режим узла $n_j$.

После выбора $\kappa_t$ вычисляется экспонируемая ёмкость каждого узла,
затем аукционный слой распределяет задачи в рамках выбранных ограничений.
Иначе говоря, действие обучаемого агента не является прямым назначением задач,
а определяет условия, в которых проводится аукционный раунд.

\section{Глобальное состояние системы}

Глобальное состояние $s_t \in \mathcal{S}$ описывает полную информацию,
необходимую для моделирования дальнейшей эволюции системы. Оно используется
на этапе централизованного обучения, но не предполагается доступным каждому
агенту во время исполнения.

В базовой модели состояние имеет вид:
\begin{equation}
	s_t =
	\bigl(
	G(t),
	\{L_j(t)\}_{j=1}^{k},
	\{Q_j(t)\}_{j=1}^{k},
	\{f_j(t)\}_{j=1}^{k},
	\mathcal{T}(t),
	\{w_{ij}(t)\}_{i,j},
	z_t
	\bigr),
	\label{eq:ch4_state}
\end{equation}
где:
\begin{itemize}
	\item $G(t)$ --- граф текущей инфраструктуры;
	\item $L_j(t)$ --- загрузка периферийного узла $n_j$;
	\item $Q_j(t)$ --- очередь задач узла $n_j$;
	\item $f_j(t)$ --- индикатор отказа или недоступности узла;
	\item $\mathcal{T}(t)$ --- множество задач текущего раунда;
	\item $w_{ij}(t)$ --- сетевые веса между роботами и узлами;
	\item $z_t$ --- дополнительные агрегированные признаки среды.
\end{itemize}

К дополнительным признакам $z_t$ могут относиться статистика предыдущих
назначений, доля задач с нарушенными сроками, уровень дисбаланса загрузки,
доля отклонённых задач, агрегированные платежи и показатели дефицита ресурсов.
Эти признаки не являются самостоятельным управляющим механизмом, но помогают
обучаемой политике учитывать последствия предыдущих решений.

\section{Локальные наблюдения периферийных узлов}

Во время исполнения агент $a_j$ не имеет доступа к полному состоянию $s_t$.
Он получает локальное наблюдение
\begin{equation}
	o^j_t = \Omega_j(s_t),
	\label{eq:ch4_obs}
\end{equation}
которое включает только доступную или агрегированную информацию.

Вектор локального наблюдения может содержать:
\begin{itemize}
	\item долю свободных вычислительных ресурсов узла $n_j$;
	\item долю свободной памяти и ускорительных ресурсов;
	\item длину локальной очереди задач;
	\item распределение задач в очереди по классам критичности;
	\item оценку среднего времени ожидания в очереди;
	\item локальные оценки задержек до ближайших роботов;
	\item признак отказа или деградации узла;
	\item операционный режим, выбранный на предыдущем шаге;
	\item число задач, назначенных узлу в предыдущем аукционном раунде;
	\item средний вклад назначенных задач в общественное благосостояние;
	\item локальные VCG-показатели: платежи, компенсации или оценки внешнего
	эффекта;
	\item локальные показатели дефицита ресурсов.
\end{itemize}

Обозначим аукционные признаки, доступные узлу $n_j$, через $g^j_t$.
Тогда локальное наблюдение удобно представить как
\[
o^j_t =
\bigl(
o^{j,base}_t,
g^j_t
\bigr),
\]
где $o^{j,base}_t$ --- базовая техническая часть наблюдения, а $g^j_t$ ---
экономически интерпретируемые признаки, сформированные по итогам предыдущих
аукционных раундов.

Важно подчеркнуть, что признаки $g^j_t$ не содержат управляемых ставок и не
задают стратегию заявок роботов. Они являются агрегированными индикаторами
того, насколько дефицитным был ресурс узла, какой внешний эффект создавали
назначенные задачи и каков был вклад узла в общественное благосостояние.

\section{Локальные истории наблюдений}

Поскольку система частично наблюдаема, одного текущего наблюдения может быть
недостаточно для корректного выбора режима. Например, одинаковая текущая
загрузка узла может означать разные ситуации: устойчивую нормальную работу,
быстро растущую очередь или восстановление после пика нагрузки. Поэтому
агенту требуется учитывать историю.

Локальная история агента $a_j$ определяется как
\begin{equation}
	h^j_t =
	\left(
	o^j_0,\kappa_j(0),r_0,
	o^j_1,\kappa_j(1),r_1,
	\ldots,
	o^j_t
	\right).
	\label{eq:ch4_local_history}
\end{equation}

Политика агента имеет вид
\begin{equation}
	\pi_j(\kappa_j \mid h^j_t),
	\label{eq:ch4_policy}
\end{equation}
то есть выбирает операционный режим на основе локальной истории наблюдений.
На практике такая зависимость может аппроксимироваться рекуррентной
нейронной сетью. Это согласуется с архитектурой QMIX, где локальные сети
агентов часто используют рекуррентный слой для работы в частично наблюдаемой
среде \cite{Rashid2018QMIX}.

\section{Пространство действий: операционные режимы узла}

Действием агента $a_j$ является выбор операционного режима периферийного
узла:
\begin{equation}
	\kappa_j(t) \in \mathcal{K}_j.
	\label{eq:ch4_action}
\end{equation}

В базовой дискретной постановке множество режимов можно задать как:
\begin{equation}
	\mathcal{K}_j =
	\{
	\kappa^{safe},
	\kappa^{normal},
	\kappa^{aggressive},
	\kappa^{reserve},
	\kappa^{critical}
	\}.
	\label{eq:ch4_actions}
\end{equation}

Содержательно эти режимы интерпретируются следующим образом:
\begin{enumerate}
	\item $\kappa^{safe}$ --- защитный режим: узел минимизирует приём новых
	задач, если очередь велика или существует риск нарушения сроков;
	\item $\kappa^{normal}$ --- нормальный режим: узел экспонирует стандартную
	долю свободной ёмкости;
	\item $\kappa^{aggressive}$ --- агрессивный режим: узел экспонирует большую
	долю свободной ёмкости при благоприятном состоянии;
	\item $\kappa^{reserve}$ --- режим резервирования: часть ёмкости сохраняется
	для возможных будущих критических задач;
	\item $\kappa^{critical}$ --- режим критического обслуживания: узел принимает
	прежде всего задачи классов, связанных с безопасностью, локализацией или
	срочным перепланированием.
\end{enumerate}

Каждый режим задаёт коэффициенты экспонирования ресурсов:
\[
\eta^r_j(\kappa_j)\in[0,1],
\quad
r \in \{cpu,mem,gpu\},
\]
а также, при необходимости, поправки к приоритетам классов задач:
\[
\chi_{j,c}(\kappa_j),
\]
где $c$ --- класс задачи.

Тогда экспонируемая ёмкость узла для ресурса $r$ равна
\begin{equation}
	\bar{A}^{r}_j(t,\kappa_j)
	=
	\eta^r_j(\kappa_j)
	\left(
	C^r_j - L^r_j(t)
	\right).
	\label{eq:ch4_exp_capacity}
\end{equation}

Если узел находится в отказе или физически не способен выполнить режим,
соответствующее действие исключается из множества допустимых. Это ограничение
является физическим и ресурсным, а не аукционным сигналом. В программной
реализации оно может задаваться маской допустимых действий, но в математической
модели оно трактуется как ограничение среды.

\section{Совместная политика и целевая функция обучения}

Пусть
\begin{equation}
	\pi = (\pi_1,\ldots,\pi_k)
	\label{eq:ch4_global_policy}
\end{equation}
--- совместная политика \eqref{eq:ch4_policy} всех периферийных узлов.

Цель обучения состоит в максимизации ожидаемой дисконтированной суммы командных вознаграждений:
\begin{equation}
	J(\pi)
	=
	\mathbb{E}_{\pi}
	\left[
	\sum_{t=0}^{\infty}
	\gamma^t r_t
	\right].
	\label{eq:ch4_target_learning}
\end{equation}

Оптимальная политика задаётся как
\begin{equation}
	\pi^* \in \arg\max_{\pi} J(\pi).
	\label{eq:ch4_optimal_policy}
\end{equation}
Вознаграждение является командным, то есть одинаковым для всех агентов. Это
соответствует кооперативной постановке: периферийные узлы принадлежат одной
инфраструктуре и должны оптимизировать глобальное качество обслуживания AMR,
а не только локальные показатели отдельных узлов.

\section{Функция вознаграждения}

Функция вознаграждения должна согласовывать обучение с целями системы~\cite{Yao2025AdaptiveOrchestration}:
обеспечивать высокое общественное благосостояние, снижать нарушения сроков,
ограничивать отклонение задач, предотвращать перегрузку и сохранять
сбалансированное использование ресурсов.

Базовое вознаграждение задаётся как:
\begin{equation}
	r_t =
	\alpha_{sw} \widetilde{SW}(t)
	-
	\alpha_{viol} \mathrm{Viol}(t)
	-
	\alpha_{drop} \mathrm{Drop}(t)
	-
	\alpha_{imb} \mathrm{Imb}(t)
	-
	\alpha_{over} \mathrm{Over}(t),
	\label{eq:ch4_reward_base}
\end{equation}
где:
\begin{itemize}
	\item $\widetilde{SW}(t)$ --- нормированное общественное благосостояние
	аукционного распределения;
	\item $\mathrm{Viol}(t)$ --- доля задач с нарушенными сроками;
	\item $\mathrm{Drop}(t)$ --- доля отклонённых или нераспределённых задач;
	\item $\mathrm{Imb}(t)$ --- показатель дисбаланса загрузки между узлами;
	\item $\mathrm{Over}(t)$ --- штраф за перегрузку узлов;
	\item $\alpha_{sw},\alpha_{viol},\alpha_{drop},\alpha_{imb},\alpha_{over}$
	--- неотрицательные коэффициенты.
\end{itemize}


Показатель дисбаланса может быть задан как дисперсия загрузки узлов:
\[
\mathrm{Imb}(t)
=
\frac{1}{k}
\sum_{j=1}^{k}
\left(
u_j(t) - \bar{u}(t)
\right)^2,
\label{eq:ch4_imb}
\]
где $u_j(t)$ --- агрегированная загрузка узла, а $\bar{u}(t)$ --- средняя
загрузка по системе.

Штраф за перегрузку:
\[
\mathrm{Over}(t)
=
\sum_{j=1}^{k}
\left[
u_j(t)-u^{max}
\right]_+,
\]
где $u^{max}$ --- допустимый уровень загрузки.

Для включения аукционных сигналов можно использовать дополнительный компонент:
\begin{equation}
	r^{auc}_t =
	\alpha_p \widetilde{P}(t)
	+
	\alpha_q \widetilde{Q}(t)
	+
	\alpha_g \widetilde{E}(t),
	\label{eq:ch4_reward_auc}
\end{equation}
где $\widetilde{P}(t)$ --- нормированный показатель платежей, $\widetilde{Q}(t)$
--- нормированный показатель внутренних компенсаций узлов, а
$\widetilde{E}(t)$ --- нормированная оценка внешних эффектов распределения.

Итоговая функция вознаграждения:
\[
\tilde{r}_t =
(1-\lambda_{auc}) r_t
+
\lambda_{auc} r^{auc}_t,
\quad
\lambda_{auc}\in[0,1].
\]

При этом аукционная компонента не должна превращать обучение в оптимизацию
платежей ради платежей. Платежи используются как индикаторы внешнего эффекта
и дефицита ресурсов, а основной критерий остаётся связанным с качеством
обслуживания и общественным благосостоянием.

\section{Централизованное обучение и децентрализованное исполнение}

Для рассматриваемой задачи естественной является парадигма централизованного
обучения и децентрализованного исполнения (Centralized Training with
Decentralized Execution, CTDE)~\cite{Amato2024CTDE}. На этапе обучения допускается использовать глобальное состояние системы, полные журналы назначений, совместные действия всех агентов и командное вознаграждение. На этапе исполнения каждый
периферийный узел использует только собственную локальную историю наблюдений.

В AMR-системах централизованное обучение может выполняться:
\begin{enumerate}
	\item в облачном контуре;
	\item на сервере управления роботизированной системой;
	\item в цифровом двойнике рабочей среды;
	\item по журналам работы реальной системы;
	\item в имитационной среде, воспроизводящей поток задач, движение роботов
	и состояние периферийных узлов.
\end{enumerate}

После обучения параметры локальных политик передаются на периферийные узлы.
В режиме исполнения узел $n_j$ выполняет следующие действия:
\begin{enumerate}
	\item получает локальное наблюдение $o^j_t$;
	\item обновляет внутреннее состояние локальной модели;
	\item выбирает операционный режим $\kappa_j(t)$;
	\item экспонирует соответствующую долю ресурсов;
	\item участвует в последующем аукционном распределении задач.
\end{enumerate}

Такой подход снижает требования к постоянной централизованной координации.
Глобальная информация используется на этапе обучения, но не требуется для
каждого локального решения во время работы системы.

\section{Обоснование выбора QMIX}

Для реализации кооперативного многоагентного управления в работе используется
алгоритм QMIX \cite{Rashid2018QMIX}. Он относится к методам факторизации
функции ценности и предназначен для задач, где агенты имеют локальные
наблюдения, но обучаются на общем командном вознаграждении.

QMIX строит индивидуальные функции ценности агентов:
\[
Q_j(h^j_t,\kappa_j;\theta_j),
\]
а затем объединяет их в совместную функцию ценности:
\[
Q_{tot}(s_t,\kappa_t)
=
f_{mix}
\left(
Q_1,\ldots,Q_k,s_t;\theta_{mix}
\right).
\]

Ключевым ограничением является монотонность:
\[
\frac{\partial Q_{tot}}{\partial Q_j} \geq 0,
\quad
\forall j.
\]
Это условие обеспечивает согласованность между локальным выбором режимов и
централизованно обученной совместной функцией ценности. Благодаря этому во
время исполнения каждый агент может выбирать режим локально:
\[
\kappa_j^*(t)
=
\arg\max_{\kappa_j \in \mathcal{K}_j}
Q_j(h^j_t,\kappa_j;\theta_j).
\]

Для рассматриваемой постановки QMIX подходит по следующим причинам:
\begin{enumerate}
	\item действия периферийных узлов дискретны и соответствуют конечному
	набору операционных режимов;
	\item требуется кооперативная оптимизация общего качества обслуживания;
	\item исполнение должно быть децентрализованным;
	\item обучение может использовать глобальное состояние и командное
	вознаграждение;
	\item метод избегает полного перебора пространства совместных действий.
\end{enumerate}

\section{Архитектура локальных сетей агентов}

Каждый агент $a_j$ аппроксимирует локальную функцию ценности
\begin{equation}
	Q_j(h^j_t,\kappa_j;\theta_j),
	\label{eq:ch4_Q_local}
\end{equation}
которая оценивает качество выбора операционного режима узла в текущей ситуации.

Поскольку наблюдения частичны и зависят от истории, локальная сеть (agent network) может
содержать рекуррентный слой:
\[
h^{net,j}_t =
\mathrm{GRU}_j
\left(
h^{net,j}_{t-1},
[o^j_t,\kappa_j(t-1)]
\right),
\]
где $h^{net,j}_t$ --- скрытое состояние сети агента.

Значения функции ценности для всех режимов вычисляются как
\[
Q_j(h^j_t,\cdot)
=
f_j(h^{net,j}_t;\theta_j).
\]

Вход локальной сети включает технические признаки узла и агрегированные
аукционные признаки.

\section{Сеть смешивания}

Сеть смешивания (mixing network) объединяет локальные значения $Q_j$ в совместную функцию:
\begin{equation}
	Q_{tot}
	=
	f_{mix}
	\left(
	Q_1,\ldots,Q_k,s_t;\theta_{mix}
	\right).
	\label{eq:ch4_Q_total}
\end{equation}

Параметры сети смешивания могут зависеть от глобального состояния $s_t$ через
гиперсети. Это позволяет учитывать общий контекст системы: число активных
задач, распределение загрузки, сетевую обстановку и аукционные показатели
предыдущих раундов.

Монотонность обеспечивается неотрицательностью весов сети смешивания:
\[
\frac{\partial Q_{tot}}{\partial Q_j} \geq 0.
\]
На практике это реализуется ограничениями на веса или использованием
неотрицательных преобразований, как в исходной архитектуре QMIX
\cite{Rashid2018QMIX}.

\section{Правило обучения}

Опыт, получаемый системой, сохраняется в буфере в виде переходов:
\begin{equation}
	e_t =
	\left(
	s_t,
	\{o^j_t\}_{j=1}^{k},
	\kappa_t,
	\tilde{r}_t,
	s_{t+1},
	\{o^j_{t+1}\}_{j=1}^{k},
	done_t
	\right).
	\label{eq:ch4_replay_buffer_vector}
\end{equation}

Для обновления параметров используется временно-разностная цель (TD-target):
\begin{equation}
	y_t =
	\tilde{r}_t
	+
	\gamma(1-done_t)
	\max_{\kappa'_{t+1}}
	Q^{-}_{tot}(s_{t+1},\kappa'_{t+1}),
	\label{eq:ch4_Q_target}
\end{equation}
где $Q^{-}_{tot}$ --- целевая сеть (target network).

С учётом монотонности максимум по совместному действию может быть получен
через локальные greedy-выборы:
\[
\kappa_j' =
\arg\max_{\kappa_j \in \mathcal{K}_j}
Q^{-}_j(h^j_{t+1},\kappa_j).
\]

Функция потерь:
\begin{equation}
	\mathcal{L}(\theta)
	=
	\mathbb{E}_{e_t \sim \mathcal{B}}
	\left[
	\left(
	y_t -
	Q_{tot}(s_t,\kappa_t;\theta)
	\right)^2
	\right],
	\label{eq:ch4_loss}
\end{equation}
где $\mathcal{B}$ --- буфер опыта (experience replay buffer).

Обновление параметров выполняется на этапе централизованного обучения. В
режиме исполнения градиентные обновления не требуются: каждый узел использует
уже обученную локальную сеть для выбора режима.

\section{Связь с аукционным слоем}

Связь между обучаемым и аукционным контурами имеет направленный характер.

Сначала обучаемый контур выбирает режимы:
\[
\kappa_t =
(\kappa_1(t),\ldots,\kappa_k(t)).
\]

Затем для каждого узла вычисляется экспонируемая ёмкость:
\[
\bar{A}_j(t,\kappa_j).
\]

После этого аукционный слой решает задачу:
\[
x^*(t)
\in
\arg\max_{x \in \mathcal{X}(t,\kappa(t))}
SW(x,t).
\]

Результаты аукциона влияют на следующие наблюдения и вознаграждение:
\[
g^j_{t+1}
=
\Phi_j
\left(
x^*(t),
p^R(t),
q^N(t),
SW(t),
\sigma_j(t)
\right),
\]
где $\Phi_j$ --- оператор локальной агрегации аукционных признаков для узла
$n_j$.

Таким образом, аукцион не является внешним постобработчиком решений QMIX, а
встроен в динамику среды: режимы узлов задают условия аукциона, а результаты
аукциона формируют часть последующего опыта для обучения.

\section{Псевдокод метода адаптивного управления режимами узлов}

Ниже приведен один шаг взаимодействия обучаемого контура со средой (см. рис.~\ref{img:arch-qmix}).

\textbf{Алгоритм 4.1. Метод адаптивного управления режимами узлов}

\begin{enumerate}
	\item каждый периферийный узел получает локальное наблюдение $o^j_t$;
	\item локальная сеть агента выбирает операционный режим $\kappa_j(t)$;
	\item выбранные режимы определяют экспонируемую ёмкость узлов;
	\item аукционный слой распределяет задачи текущего раунда;
	\item назначенные задачи поступают в очереди или на исполнение;
	\item среда обновляет загрузку, очереди, задержки и отказовые состояния;
	\item рассчитываются общественное благосостояние, штрафы и аукционные
	показатели;
	\item формируется командное вознаграждение $\tilde{r}_t$;
	\item переход сохраняется в буфер опыта;
	\item на этапе обучения выполняется обновление параметров QMIX.
\end{enumerate}

\begin{figure}[h]
	\centering
	\includegraphics[width=0.8\textwidth]{arch-qmix}
	\caption{Концептуальная архитектура алгоритма}
	\label{img:arch-qmix}
\end{figure}

\section{Ограничения обучаемого контура}

Предложенный обучаемый контур имеет ряд ограничений.

Во-первых, QMIX не гарантирует нахождение глобально оптимальной политики.
Он является приближённым методом обучения в сложной частично наблюдаемой
многоагентной среде.

Во-вторых, ограничение монотонности может быть слишком сильным для сценариев,
где вклад локального действия одного узла в совместную ценность зависит от
действий других узлов немонотонным образом. В таких случаях могут потребоваться
более выразительные методы факторизации функции ценности.

В-третьих, качество обученной политики зависит от репрезентативности
обучающих сценариев. Если реальные потоки задач, задержки, отказы или
поведение роботов существенно отличаются от симуляционной среды, требуется
дообучение или перенос политики.

В-четвёртых, обучаемый контур не заменяет аукционный механизм и не устраняет
его вычислительную сложность. Он управляет режимами узлов, но сам аукционный
раунд всё ещё должен быть решён в рамках выбранных ограничений.

В-пятых, использование аукционных признаков в наблюдениях и вознаграждении
не означает, что вся система наследует свойства VCG. Стратегическая
совместимость относится только к локальному аукционному раунду при стандартных
предпосылках, как было показано в предыдущей главе.

\section{Выводы по главе}

В текущей главе построена модель адаптивного многоагентного управления
периферийными узлами в системах автономных мобильных роботов с поддержкой
периферийных вычислений. Показано, что раундовый аукционный механизм
необходимо дополнить обучаемой политикой, поскольку он не решает задачу
долгосрочной адаптации к изменению очередей, задержек, отказов и входного
потока задач.

Задача выбора операционных режимов периферийных узлов формализована как
децентрализованный частично наблюдаемый марковский процесс принятия решений.
Обучаемыми агентами являются периферийные узлы, а их действиями --- режимы
допуска, резервирования и обслуживания задач. Такая постановка сохраняет
разделение ролей: роботы формируют задачи и заявки, а периферийные узлы
адаптивно управляют доступной ёмкостью.

Обосновано применение парадигмы централизованного обучения и
децентрализованного исполнения. Централизованное обучение может выполняться в
облаке, серверном контуре или цифровом двойнике, а исполнение политики
происходит локально на периферийных узлах.

В качестве базового алгоритма выбран QMIX, поскольку он согласуется с
дискретным пространством режимов, командным вознаграждением,
кооперативной постановкой и требованием децентрализованного исполнения.
Показана связь обучаемого контура с аукционным слоем: политика выбирает
режимы узлов, а аукцион распределяет задачи в рамках выбранных ограничений.
Полученная модель является основой для построения полного гибридного метода
AMPLE-AMR в следующей главе.

